% =============================================================================
% Serial-Protokoll.tex – Robuste Frames mit COBS + CRC
% =============================================================================
% Fragment für \input{} – keine Präambel, nutzt techdoc.cls
% =============================================================================

\section{Serial-Protokoll -- Robuste Frames mit COBS + CRC}
\label{sec:serial-protokoll}

\subsection{Begriffsrahmen und Zielsetzung}
\label{subsec:serial-begriffsrahmen}

Der ESP32-S3 kommuniziert mit dem Raspberry Pi über USB-CDC (virtueller Serial-Port).
Die Herausforderung: Serielle Kommunikation ist ein kontinuierlicher Byte-Strom ohne
inhärente Nachrichtengrenzen.

\begin{table}[H]
\centering
\caption{Zentrale Begriffe des Serial-Protokolls}
\label{tab:serial-begriffe}
\begin{tabular}{@{}ll@{}}
\toprule
\textbf{Begriff} & \textbf{Definition} \\
\midrule
USB-CDC  & USB Communication Device Class -- emuliert einen seriellen Port über USB \\
Frame    & Ein abgegrenztes Datenpaket mit definiertem Anfang und Ende \\
COBS     & Consistent Overhead Byte Stuffing -- Framing-Verfahren mit \code{0x00} als Delimiter \\
CRC-16   & Cyclic Redundancy Check -- Prüfsumme zur Fehlererkennung \\
Payload  & Nutzdaten innerhalb eines Frames \\
ACK      & Acknowledgment -- Empfangsbestätigung \\
\bottomrule
\end{tabular}
\end{table}

\begin{tipbox}[Ziel]
Zuverlässige, bidirektionale Kommunikation mit Fehlererkennung und Retry-Mechanik.
\end{tipbox}

% -----------------------------------------------------------------------------
\subsection{Das Grundproblem -- Byte-Strom vs.\ Nachrichten}
\label{subsec:serial-grundproblem}

Wenn du Bytes über eine serielle Leitung schickst, sieht der Empfänger nur einen
kontinuierlichen Strom:

\begin{lstlisting}[style=base,numbers=none]
Sender:   [Nachricht A: 5 Bytes][Nachricht B: 3 Bytes][Nachricht C: 7 Bytes]
Leitung:  0x12 0x34 0x56 0x78 0x9A 0xBC 0xDE 0xF0 0x11 0x22 0x33 0x44 0x55 0x66 0x77
Empfänger: "Wo fängt Nachricht B an?"
\end{lstlisting}

Ohne Framing weiß der Empfänger nicht, wo eine Nachricht endet und die nächste beginnt.
Drei klassische Lösungsansätze existieren:

\begin{table}[H]
\centering
\caption{Vergleich der Framing-Ansätze}
\label{tab:framing-ansaetze}
\begin{tabular}{@{}lll@{}}
\toprule
\textbf{Ansatz} & \textbf{Prinzip} & \textbf{Nachteil} \\
\midrule
Feste Länge & Jede Nachricht hat exakt $N$ Bytes & Unflexibel, verschwendet Bandbreite \\
Längenfeld  & Erstes Byte gibt Länge an          & Synchronisationsverlust bei Byte-Verlust \\
Delimiter   & Spezielles Zeichen markiert Ende   & Was, wenn das Zeichen in den Daten vorkommt? \\
\bottomrule
\end{tabular}
\end{table}

COBS kombiniert den Delimiter-Ansatz mit Byte-Stuffing -- das Trennzeichen \code{0x00}
kann garantiert nicht in den Daten vorkommen.

% -----------------------------------------------------------------------------
\subsection{COBS -- Consistent Overhead Byte Stuffing}
\label{subsec:cobs}

Wir transformieren die Daten so, dass \code{0x00} niemals im Inhalt auftritt.
Dann können wir \code{0x00} als eindeutigen Delimiter verwenden.

\paragraph{Wie COBS funktioniert}
COBS ersetzt jede \code{0x00} in den Originaldaten durch einen \glqq Zeiger\grqq\ auf
die nächste Null (oder das Paketende):

\begin{lstlisting}[style=base,numbers=none]
Original:     [0x11] [0x22] [0x00] [0x33] [0x00] [0x44]
                             ^             ^
                          Nullen müssen weg

COBS-kodiert: [0x03] [0x11] [0x22] [0x02] [0x33] [0x02] [0x44]
               |                    |             |
               +-- "nach 3 Bytes    +-- "nach 2   +-- "nach 2 Bytes
                   kommt Ersatz"        Bytes"        kommt Ende"
\end{lstlisting}

Der Empfänger liest den ersten Byte-Wert (hier \code{0x03}), kopiert die folgenden
2~Bytes, fügt eine \code{0x00} ein, liest den nächsten Zeiger, und so weiter.

\paragraph{Wire-Frame-Aufbau}

\begin{lstlisting}[style=base,numbers=none]
+---------------------------------------------+-----------+
|  COBS-kodierte Daten (keine 0x00 enthalten) |   0x00    |
|                                             | Delimiter |
+---------------------------------------------+-----------+
\end{lstlisting}

COBS fügt im Worst Case \SI{1}{\byte} pro \SI{254}{\byte} Originaldaten hinzu.
Bei typischen Paketgrößen unter \SI{100}{\byte} ist der Overhead vernachlässigbar
(1--2~Bytes).

\begin{infobox}[Konsequenz für den Empfänger]
Du sammelst Bytes, bis du \code{0x00} siehst. Dann dekodierst du alles davor mit COBS
und hast das Original-Paket.
\end{infobox}

% -----------------------------------------------------------------------------
\subsection{CRC-16 -- Integritätsprüfung}
\label{subsec:crc16}

Auch mit korrektem Framing können Bytes auf der Leitung verfälscht werden -- durch EMI,
Timing-Probleme oder Buffer-Überläufe. Der Sender berechnet eine Prüfsumme über die
Daten und hängt sie an. Der Empfänger berechnet dieselbe Prüfsumme und vergleicht.

\paragraph{CRC-16/CCITT-FALSE Parameter}

\begin{table}[H]
\centering
\caption{CRC-16/CCITT-FALSE Konfiguration}
\label{tab:crc16-parameter}
\begin{tabular}{@{}lll@{}}
\toprule
\textbf{Parameter} & \textbf{Wert} & \textbf{Bedeutung} \\
\midrule
Polynom      & \code{0x1021} & Mathematische Basis der Berechnung \\
Init         & \code{0xFFFF} & Startwert \\
XorOut       & \code{0x0000} & Finaler XOR (hier: keiner) \\
RefIn/RefOut & \code{false}  & Keine Bit-Spiegelung \\
\bottomrule
\end{tabular}
\end{table}

Dieses Polynom erkennt alle Einzelbitfehler, alle Doppelbitfehler, und die meisten
Burst-Fehler bis \SI{16}{\bit} Länge. Die Erkennungswahrscheinlichkeit für zufällige
Fehler liegt bei $1 - 2^{-16} \approx 99{,}998\,\%$.

\paragraph{Berechnung (Python-Referenz)}

\begin{lstlisting}[style=python,caption={CRC-16/CCITT-FALSE Implementierung},label={lst:crc16-python}]
def crc16_ccitt_false(data: bytes) -> int:
    crc = 0xFFFF
    for b in data:
        crc ^= (b << 8) & 0xFFFF
        for _ in range(8):
            if crc & 0x8000:
                crc = ((crc << 1) ^ 0x1021) & 0xFFFF
            else:
                crc = (crc << 1) & 0xFFFF
    return crc
\end{lstlisting}

Wenn CRC nicht stimmt, verwirft der Empfänger das Paket. Der Sender muss es erneut
senden (siehe \cref{subsec:ack-retry}).

% -----------------------------------------------------------------------------
\subsection{Packet-Layout -- Der Header}
\label{subsec:packet-layout}

Bevor COBS angewendet wird, bauen wir das \glqq Raw-Packet\grqq\ mit einem strukturierten Header:

\begin{table}[H]
\centering
\caption{Packet-Header-Struktur}
\label{tab:packet-header}
\begin{tabular}{@{}cllp{6cm}@{}}
\toprule
\textbf{Offset} & \textbf{Feld} & \textbf{Größe} & \textbf{Beschreibung} \\
\midrule
0       & magic0     & 1 & \code{'S'} (\code{0x53}) \\
1       & magic1     & 1 & \code{'P'} (\code{0x50}) \\
2       & ver\_major & 1 & 1 \\
3       & ver\_minor & 1 & 0 \\
4       & type       & 1 & Message-Typ \\
5       & flags      & 1 & Bitflags \\
6--7    & seq        & 2 & Sequenznummer (Little-Endian) \\
8--9    & len        & 2 & Payload-Länge (Little-Endian) \\
10..    & payload    & len & Nutzdaten \\
10+len  & crc16      & 2 & CRC über Bytes 0..(10+len-1) \\
\bottomrule
\end{tabular}
\end{table}

\paragraph{Warum diese Felder?}

\textbf{Magic Bytes (\code{SP}):} Wenn wir mitten in einem Datenstrom einsteigen oder
ein Frame korrupt ist, erkennen wir sofort, dass etwas nicht stimmt. Die Wahrscheinlichkeit,
dass zufällige Daten mit \code{0x53 0x50} beginnen, liegt bei $1/65536 \approx 0{,}0015\,\%$.

\textbf{Version:} Ermöglicht Protokoll-Evolution. Regel: Major-Mismatch führt zum Verwerfen
des Frames. Minor-Mismatch wird akzeptiert (abwärtskompatibel).

\textbf{Sequenznummer:} Identifiziert Pakete für ACK-Zuordnung. Bereich 1--65535, Wert~0
wird übersprungen.

\paragraph{Flags}

\begin{table}[H]
\centering
\caption{Flag-Bits im Header}
\label{tab:header-flags}
\begin{tabular}{@{}cll@{}}
\toprule
\textbf{Bit} & \textbf{Name} & \textbf{Bedeutung} \\
\midrule
0 & \code{ACK\_REQ} & Sender erwartet Bestätigung \\
1 & \code{IS\_ACK}  & Dieses Paket ist eine Bestätigung \\
2 & \code{IS\_ERR}  & Dieses Paket meldet einen Fehler \\
\bottomrule
\end{tabular}
\end{table}

% -----------------------------------------------------------------------------
\subsection{Sende- und Empfangsprozess}
\label{subsec:sende-empfang}

\paragraph{Senden (4 Schritte)}

\begin{enumerate}
\item \textbf{Header + Payload zusammenbauen:} Header (\SI{10}{\byte}) gefolgt von Payload.
\item \textbf{CRC-16 berechnen und anhängen:} CRC über gesamten Inhalt, dann 2~Bytes anfügen.
\item \textbf{COBS-encode:} Eliminiert alle \code{0x00} aus dem Datenstrom.
\item \textbf{Delimiter anhängen und senden:} \code{0x00} markiert das Frame-Ende.
\end{enumerate}

\paragraph{Empfangen (4 Schritte)}

\begin{enumerate}
\item Bytes sammeln bis \code{0x00}.
\item COBS-decode.
\item CRC prüfen (letzte 2~Bytes vs.\ Berechnung).
\item Header parsen, Payload extrahieren.
\end{enumerate}

% -----------------------------------------------------------------------------
\subsection{Message-Typen}
\label{subsec:message-typen}

\paragraph{ESP32 → Pi (Events)}

\begin{table}[H]
\centering
\caption{Event-Typen vom ESP32 zum Pi}
\label{tab:msg-esp-to-pi}
\begin{tabular}{@{}llp{7cm}@{}}
\toprule
\textbf{Type} & \textbf{Hex} & \textbf{Zweck} \\
\midrule
HELLO           & \code{0x01} & Fähigkeiten nach Reset melden \\
BTN\_EVENT      & \code{0x10} & Button gedrückt/losgelassen \\
SELECTION\_STATE & \code{0x12} & Aktuell gewählter Kanal \\
SNAPSHOT        & \code{0x11} & Kompletter Button/LED-Zustand \\
\bottomrule
\end{tabular}
\end{table}

\paragraph{Pi → ESP32 (Commands)}

\begin{table}[H]
\centering
\caption{Command-Typen vom Pi zum ESP32}
\label{tab:msg-pi-to-esp}
\begin{tabular}{@{}llp{7cm}@{}}
\toprule
\textbf{Type} & \textbf{Hex} & \textbf{Zweck} \\
\midrule
HELLO\_ACK         & \code{0x02} & HELLO bestätigen \\
CMD\_SET\_SELECTION & \code{0x20} & Kanal vom Dashboard setzen \\
CMD\_REQ\_SNAPSHOT  & \code{0x23} & Debug-Snapshot anfordern \\
\bottomrule
\end{tabular}
\end{table}

\paragraph{Bidirektional (Control)}

\begin{table}[H]
\centering
\caption{Bidirektionale Control-Typen}
\label{tab:msg-bidirektional}
\begin{tabular}{@{}llp{7cm}@{}}
\toprule
\textbf{Type} & \textbf{Hex} & \textbf{Zweck} \\
\midrule
ACK  & \code{0x30} & Empfangsbestätigung \\
ERR  & \code{0x50} & Fehler melden \\
PING & \code{0x03} & Keepalive-Anfrage \\
PONG & \code{0x04} & Keepalive-Antwort \\
\bottomrule
\end{tabular}
\end{table}

\paragraph{SELECTION\_STATE -- Das zentrale Event}

Dieses Paket synchronisiert den Zustand zwischen ESP32, Pi und Browser:

\begin{table}[H]
\centering
\caption{SELECTION\_STATE Payload-Struktur}
\label{tab:selection-state}
\begin{tabular}{@{}cllp{5cm}@{}}
\toprule
\textbf{Offset} & \textbf{Feld} & \textbf{Typ} & \textbf{Beschreibung} \\
\midrule
0--1 & selected\_id & u16 & 0 = keiner, 1..100 = Kanal \\
2    & source       & u8  & Wer hat ausgelöst? \\
3    & reserved     & u8  & 0 \\
4--7 & t\_ms        & u32 & \code{millis()} seit Boot \\
\bottomrule
\end{tabular}
\end{table}

\begin{table}[H]
\centering
\caption{Source-Werte für SELECTION\_STATE}
\label{tab:source-werte}
\begin{tabular}{@{}clp{7cm}@{}}
\toprule
\textbf{Wert} & \textbf{Name} & \textbf{Bedeutung} \\
\midrule
0 & \code{local\_btn} & Physischer Button am Panel \\
1 & \code{host\_cmd}  & Dashboard-Klick (via Pi) \\
2 & \code{reset}      & Nach ESP32-Neustart \\
3 & \code{recover}    & Wiederherstellung nach Verbindungsabbruch \\
\bottomrule
\end{tabular}
\end{table}

% -----------------------------------------------------------------------------
\subsection{ACK/Retry-Mechanik}
\label{subsec:ack-retry}

Bei Commands wie \code{CMD\_SET\_SELECTION} muss der Pi wissen, ob der ESP32 das Kommando
erhalten und verarbeitet hat. Die Lösung: Request-Response mit Timeout und Retry.

\paragraph{Ablauf}

\begin{lstlisting}[style=base,numbers=none]
Pi                                 ESP32
 |                                   |
 |--- CMD_SET_SELECTION ------------>|  seq=42, ACK_REQ=1
 |         (Timer startet)           |
 |                                   |  (verarbeitet...)
 |<---------- ACK -------------------|  seq=42, IS_ACK=1
 |         (Timer stoppt)            |
 |                                   |
 |<-- SELECTION_STATE ---------------|  selected_id=7, source=host_cmd
\end{lstlisting}

\paragraph{Fehlerfall}

\begin{lstlisting}[style=base,numbers=none]
Pi                                 ESP32
 |                                   |
 |--- CMD_SET_SELECTION ------------>|  seq=43, ACK_REQ=1
 |         (Timer: 600ms)            |
 |              ...                  |  (keine Antwort)
 |         (Timeout!)                |
 |--- CMD_SET_SELECTION ------------>|  seq=43, Retry #1
 |              ...                  |
 |         (Timeout!)                |
 |--- CMD_SET_SELECTION ------------>|  seq=43, Retry #2
 |              ...                  |
 |<---------- ACK -------------------|  (endlich!)
\end{lstlisting}

\paragraph{Parameter (bridge.py)}

\begin{table}[H]
\centering
\caption{ACK/Retry-Parameter}
\label{tab:ack-retry-parameter}
\begin{tabular}{@{}llp{6cm}@{}}
\toprule
\textbf{Parameter} & \textbf{Default} & \textbf{Bedeutung} \\
\midrule
\code{--ack-timeout} & \SI{0,6}{\second} & Wartezeit auf ACK \\
\code{--retries}     & 4                 & Maximale Wiederholungen \\
\bottomrule
\end{tabular}
\end{table}

\begin{warnbox}[Fehlerbehandlung]
Nach 5 fehlgeschlagenen Versuchen (1 + 4 Retries) meldet die Bridge einen Fehler
an den Browser (\code{error code=21}).
\end{warnbox}

% -----------------------------------------------------------------------------
\subsection{Datenfluss -- Vom Button zum Browser}
\label{subsec:datenfluss}

Wenn du Button~7 am Panel drückst, passiert folgendes:

\begin{enumerate}
\item \textbf{ESP32:} CD4021BE liest Button-Matrix (alle \SI{10}{\milli\second}).
      Debounce erkennt stabile Flanke nach \SI{50}{\milli\second}.

\item \textbf{ESP32:} Baut \code{SELECTION\_STATE}-Paket.
      Header: \code{magic=SP}, \code{type=0x12}, \code{seq=1234}, \code{len=8}.
      Payload: \code{selected\_id=7}, \code{source=0}, \code{t\_ms=45678}.
      CRC16 anhängen, COBS-encode, \code{0x00} Delimiter.

\item \textbf{Pi (bridge.py):} Empfängt Frame über USB-CDC Serial.
      COBS-decode, CRC prüfen, Header parsen. Erkennt \code{SELECTION\_STATE}.

\item \textbf{Pi:} Wandelt in JSON für WebSocket:
\begin{lstlisting}[style=json]
{"type":"selection","selected":7,"source":"local_btn","tMs":45678}
\end{lstlisting}

\item \textbf{Browser (React):} \code{useBridgeWs} empfängt JSON.
      State-Update führt dazu, dass LED~7 im Dashboard leuchtet.
\end{enumerate}

\begin{infobox}[Latenz]
Vom Button-Druck bis zur Browser-Anzeige vergehen typisch \SIrange{20}{50}{\milli\second}
(abhängig von Debounce-Zeit und Netzwerk).
\end{infobox}

% -----------------------------------------------------------------------------
\subsection{Größenrahmen und Limits}
\label{subsec:groessenrahmen}

\begin{table}[H]
\centering
\caption{Protokoll-Konstanten und Limits}
\label{tab:protokoll-limits}
\begin{tabular}{@{}llp{6cm}@{}}
\toprule
\textbf{Konstante} & \textbf{Wert} & \textbf{Bedeutung} \\
\midrule
\code{HEADER\_LEN}     & \SI{10}{\byte}    & Feste Header-Größe \\
\code{CRC\_LEN}        & \SI{2}{\byte}     & CRC-16 \\
\code{MAX\_PAYLOAD}    & \SI{256}{\byte}   & Firmware-Limit \\
\code{MAX\_PACKET}     & \SI{268}{\byte}   & Header + Payload + CRC \\
\code{MAX\_WIRE\_FRAME} & $\sim$\SI{270}{\byte} & Nach COBS + Delimiter \\
\bottomrule
\end{tabular}
\end{table}

Ein \code{SELECTION\_STATE}-Paket ist \SI{10}{\byte} (Header) + \SI{8}{\byte} (Payload)
+ \SI{2}{\byte} (CRC) = \SI{20}{\byte} vor COBS, etwa 21--22~Bytes on-wire.

% -----------------------------------------------------------------------------
\subsection{Fehlerbehandlung}
\label{subsec:fehlerbehandlung}

\begin{table}[H]
\centering
\caption{Fehlerarten und Reaktionen}
\label{tab:fehlerbehandlung}
\begin{tabular}{@{}lll@{}}
\toprule
\textbf{Fehlerart} & \textbf{Erkennung} & \textbf{Reaktion} \\
\midrule
CRC-Mismatch       & CRC-Prüfung                     & Frame verwerfen \\
Magic falsch       & Header-Prüfung                  & Frame verwerfen, Resync \\
Länge inkonsistent & \code{len} vs.\ tatsächliche Größe & Frame verwerfen \\
ACK-Timeout        & Timer                           & Retry, dann Error \\
Unbekannter Type   & Dispatch                        & Ignorieren (forward-compatible) \\
\bottomrule
\end{tabular}
\end{table}

\begin{tipbox}[Resynchronisation]
Da \code{0x00} garantiert nur als Delimiter vorkommt, findet der Empfänger nach einem
korrupten Frame automatisch den nächsten gültigen Frame-Start.
\end{tipbox}

% -----------------------------------------------------------------------------
\subsection{Siehe auch}
\label{subsec:serial-siehe-auch}

\begin{itemize}
\item \code{apps/bridge/docs/PROTOCOL.md} -- SSOT für Wire-Format und WebSocket-Events
\item \code{apps/bridge/bridge.py} -- Python-Implementation
\item \code{firmware/src/protocol/} -- ESP32-Implementation
\item \code{apps/dashboard/src/lib/} -- TypeScript-Typen für WebSocket-Events
\end{itemize}
